\documentclass[11pt]{article}
\usepackage[margin=1in]{geometry}
\usepackage{amsmath,amssymb,amsthm,booktabs}
\usepackage[hidelinks]{hyperref}
\usepackage{xurl}
\setlength{\emergencystretch}{3em}
\newtheorem{theorem}{Theorem}
\newtheorem{lemma}[theorem]{Lemma}
\newtheorem{corollary}[theorem]{Corollary}
\newcommand{\F}{\mathbb F}
\newcommand{\Sig}{\Sigma_2}
\title{Rational replacements of four-face twisted diagrams\\
and mixed double branched covers}
\author{Topological proof with explicit classical imports}
\date{17 November 2026}
\begin{document}
\maketitle

A four-face twisted diagram has a simple checkerboard dual with four
vertices. Arbitrary rational replacements of its twist boxes have a
strong topological restriction: an irreducible double branched cover
cannot have a mixed JSJ decomposition. This excludes all twenty-eight
retained examples whose actual covers contain a hyperbolic component
separated by an incompressible torus.

The proof works directly with branched covers of the box construction.
It makes no assumption that a nonalternating replacement diagram can be
converted to a minimal alternating diagram while preserving its
polyhedral core.

\section{The exact class}

Let $G$ be a connected, loopless, simple plane graph on four vertices.
Its medial projection has one four-valent crossing site for each edge
of $G$, with connections determined by the cyclic orders at vertices and
faces. Remove a small disk around every site and insert an arbitrary
rational two-string tangle, with the four ports glued in their prescribed
cyclic order. Let $\mathcal R_4$ be the set of knots obtained in this way.
All rational slopes, rotations and mirrors are allowed.

For a BBK twisted diagram, let $\Gamma$ be its weighted fat-vertex
graph and $G=\Gamma^*$ its simple plane dual; Section~4 of \cite{bbk}
explicitly calls this dual a simple graph. Theorem~1 of that source
gives $b=\#V(G)$. Hence every arbitrary rational replacement of the
twist boxes of a sharp-four BBK base is in $\mathcal R_4$.
We do not require the replacement to preserve reflection labels or the
Wirtinger upper bound. This is a class of replacements on a four-face
base, rather than an assertion about rank-dropping replacements on
higher-face bases.

\begin{theorem}\label{thm:class}
If $K\in\mathcal R_4$ and $\Sig(K)$ is irreducible, then $\Sig(K)$
is either hyperbolic or a graph manifold.
In particular it cannot contain an incompressible torus which cuts
off a finite-volume hyperbolic component.
\end{theorem}

\section{The proper subgraphs of \texorpdfstring{$K_4$}{K4}}

Call a two-string tangle algebraic if it is built from rational tangles
by Conway sums, products, rotations and mirrors.
The closure of such a tangle is arborescent.
We use the classical Conway--Montesinos fact that the double branched
cover of an arborescent knot is assembled from Seifert pieces along
tori, and therefore has no hyperbolic piece in its prime/JSJ
decomposition \cite[Introduction, pp.~vi--vii]{bs}.
Connected sums with two-bridge knots add lens-space prime factors and
also add no hyperbolic pieces.

Here is the local gluing used below, stated to fix the topology of the
graph reductions. Number a box's ports cyclically as
$0=\mathrm{NW},1=\mathrm{NE},2=\mathrm{SE},3=\mathrm{SW}$.
Conway addition glues port $1$ of the left box to port $0$ of the right
box, and port $2$ of the left box to port $3$ of the right box.
The remaining four ports form the new box.
Rotating either box or the whole neighborhood gives the corresponding
vertical product. In both operations two adjacent ports from each box
are joined by two untwisted arcs; algebraic tangles are closed under
these gluings.

\begin{lemma}\label{lem:proper}
If $G\ne K_4$, every knot obtained by the above rational decoration
of its medial projection has no hyperbolic piece in its double
branched cover. If that cover is irreducible, it is a graph manifold.
\end{lemma}
\begin{proof}
First suppose $G$ has a leaf. The leaf disk supplies a cap joining the
two ports at one end of the rational box on its sole incident edge.
A ball containing this box and cap meets the remaining diagram in two
points. Because the final diagram is a knot, the ball contains a single
proper arc, rather than a proper arc and a separate closed component.
Closing that arc with the complementary cap gives a numerator or
denominator closure of the rational tangle, hence a two-bridge knot
(possibly the unknot). The original knot is the connected sum of this
knot and the knot represented after deleting the leaf and its incident
edge. Thus leaf deletion only removes a two-bridge summand.

Delete leaves until none remains. A connected simple graph with at most
four vertices and minimum degree at least two is, among the proper
subgraphs under consideration, a triangle, a four-cycle, or the diamond
$K_4-e$. Indeed, with four vertices it has four or five edges:
four forces every degree to be two, while five is exactly $K_4-e$.
With three vertices the only such graph is the triangle.
If all edges disappear, only the two-bridge summands remain.

In a cycle, a degree-two vertex and the two incident edge boxes have a
regular neighborhood meeting the medial diagram in four points.
Its two internal connections are exactly the two adjacent-port gluings
just specified, up to rotations. Replace this neighborhood by the
resulting algebraic box. Repeatedly doing so leaves two vertices joined
by two parallel edges.

For the diamond, write its edge set as
\[
                 ab,\quad ac,\quad bc,\quad ad,\quad bd.
\]
The degree-two vertices are $c$ and $d$. Absorb each with its two
incident rational boxes. This leaves vertices $a,b$ joined by three
parallel edges carrying algebraic boxes. No assumption that the sum
of two rational tangles is rational is used.

In either case the remaining plane graph has two vertices and two or
three parallel edges, and no other edges or vertices.
Choose consecutive parallel edges. The region between them is an empty
bigon, so their two boxes again combine by the stated adjacent-port
gluing, up to rotation. Continue until one algebraic box remains.
The two vertex disks close its four ports in pairs, giving an
arborescent knot. This argument accounts explicitly for the parallel
edges created during degree-two reductions; it never collapses
nonadjacent edges with other graph material between them.
Together with the two-bridge leaf summands and the classical branched
cover construction, this proves the lemma.
\end{proof}

There are $38$ connected labelled simple graphs on four vertices.
One is $K_4$; the other $37$ follow the preceding cases.
A complete enumeration gives $16$ trees, $12$ triangle-with-leaf
graphs, $3$ four-cycles and $6$ diamonds.
This finite check corroborates the elementary degree argument;
it is not a substitute for the local tangle-gluing proof.

\section{The full \texorpdfstring{$K_4$}{K4} case and its actual branched cover}

The medial graph of the plane tetrahedral graph $K_4$ is the
octahedral four-valent graph, conventionally denoted $6^*$.
Let $Q$ be $S^3$ with disjoint open balls around its six crossing
sites removed. The remaining medial edges are twelve disjoint properly
embedded arcs in $Q$, lying on the projection sphere.
Each boundary sphere of $Q$ has exactly four marked arc endpoints.

\begin{lemma}\label{lem:m6}
The double cover of this punctured $6^*$ pair, branched over all
twelve arcs, is the exterior $M_6$ of the minimally twisted
six-component chain link. Every rational decoration in the $K_4$ case
therefore has double branched cover a closed Dehn filling of $M_6$.
\end{lemma}
\begin{proof}
Kolpakov--Martelli \cite[Figure~2 and Section~1.1, printed pp.~5--6]{km}
construct an involution of the six-chain-link exterior $N=M_6$.
They identify its quotient as the octahedral orbifold with underlying
space $S^3$ minus the six ideal vertices, and with all twelve
octahedral edges labelled by index two.
Their proof constructs four copies of the regular ideal octahedron
and identifies the quotient by the involution as the double of one
octahedron. The six cusps of $N$ correspond one-to-one to its vertices.

Truncate the six cusp ends in this quotient. Its underlying space is
$Q$, and its order-two singular set consists exactly of the twelve
truncated octahedral edges. Thus the quotient map is an ordinary double
branched cover of precisely the punctured $6^*$ pair described above.
At each boundary sphere it is the double cover of a sphere with four
simple branch points. Riemann--Hurwitz gives Euler characteristic
$2\cdot2-4=0$, so the connected boundary cover is a torus.
The order-two edge labels specify nontrivial branching at every marked
point; this is the four-marked-sphere cover required by rational-tangle
filling, rather than an unmarked orbifold resemblance.

The double branched cover of a rational tangle is a solid torus.
For the trivial tangle this follows by taking the double cover of a
disk branched at two points, an annulus, and then taking its product
with an interval. The general rational case is obtained by a
homeomorphism of the tangle pair, with its boundary marking retained.
The boundary cover of a four-marked sphere is uniquely determined by
sending each puncture meridian to the nontrivial element of $C_2$.
Consequently each prescribed boundary gluing lifts, and gluing the
six covering solid tori to $N$ gives a connected double cover of $S^3$
branched over the decorated knot. The two choices of a lift differ by
the deck involution and do not change the filled manifold.
By uniqueness of the double branched cover of a knot, this is
$\Sig(K)$. Attaching each solid torus is precisely Dehn filling along
a primitive slope. No restriction on these slopes is needed.
\end{proof}

The ordinary exterior of the Borromean rings, drawn by assigning
single crossings to the six sites of $6^*$, is a different space.
It is not the punctured branch pair in Lemma~\ref{lem:m6}.
In particular the use of four octahedra for $M_6$ is not an assertion
that the three-cusped Borromean exterior is six-cusped.
The primary quotient diagram and its six four-marked cusp spheres
fix the relevant construction.

\section{The exceptional fillings and the mod-two obstruction}

Martelli's classification \cite[Theorem~2.7, printed pp.~15--16]{mar}
states that every closed nonhyperbolic filling of $M_6$ is a graph
manifold, allowing the reducible cases specified in the subsequent
discussion, or one of two sporadic mixed manifolds.
Both mixed manifolds glue the figure-eight knot exterior $M_1$ to
\[
                      B=(D;(2,1),(2,1)),
\]
the orientable twisted interval bundle over the Klein bottle.
The two gluing maps are specified in that theorem.
We need no computation of their individual slopes.

\begin{lemma}\label{lem:modtwo}
Neither mixed exception can be the double branched cover of a knot.
\end{lemma}
\begin{proof}
The first homology groups of $M_1$ and $B$ are respectively
$\mathbb Z$ and $\mathbb Z\oplus\mathbb Z/2$.
The latter follows at once from the Klein-bottle group presentation
$\langle u,v\mid uvu^{-1}=v^{-1}\rangle$.
The Mayer--Vietoris map from their common boundary torus has the form
\[
 H_1(T^2;\F_2)\longrightarrow
 H_1(M_1;\F_2)\oplus H_1(B;\F_2).
\]
Its domain has dimension two and its target dimension three.
Its cokernel embeds in the first homology of the glued manifold;
in fact the connectedness terms make it the whole first homology.
Thus either mixed exception has nonzero $H_1(-;\F_2)$, for any gluing.

In contrast, every knot double branched cover has zero first
homology over $\F_2$. Here is an elementary group proof.
Take a Wirtinger presentation with meridians $m_i$, add $m_i^2=1$,
and fix a meridian $m_0$. The parity kernel in this orbifold group
is $\pi_1\Sig(K)$: take the unbranched double cover of the knot
exterior and then fill the lifts of meridian squares.
Put $a_i=m_i m_0$, so $a_0=1$.
These elements generate the kernel, because every even word is a
product of pairs $m_i m_j=a_i a_j^{-1}$.
A crossing relation becomes $m_b=m_o m_a m_o$, independently of its
original sign. Since $m_0a_i m_0=a_i^{-1}$, it rewrites as
$a_b=a_o a_a^{-1}a_o$.
In mod-two abelianization this gives $a_b=a_a$.
Following the underarcs once around the knot identifies every
$a_i$ with $a_0=0$. Thus $H_1(\Sig(K);\F_2)=0$.
The crossing-free unknot case is immediate.
This excludes both mixed exceptions.
\end{proof}

\begin{proof}[Proof of Theorem~\ref{thm:class}]
For a proper subgraph use Lemma~\ref{lem:proper}.
For $G=K_4$, Lemma~\ref{lem:m6} identifies the actual cover with a
closed filling of $M_6$. If hyperbolic, the conclusion holds.
Otherwise apply Martelli's exhaustive classification.
Lemma~\ref{lem:modtwo} excludes its mixed cases.
Irreducibility of $\Sig(K)$ removes the nontrivial connected-sum
degeneracies among the graph cases. Its remaining pieces are all
Seifert fibered, so it is a graph manifold.
\end{proof}

\section{Application and assurance}

\begin{corollary}\label{cor:targets}
All twenty-eight knots in the retained actual-cover topology
certificate lie outside $\mathcal R_4$, and hence outside arbitrary
rational replacements of the twist boxes of any sharp-four BBK base.
\end{corollary}
\begin{proof}
The retained proof \path{SEED_COVER_TOPOLOGY.tex} and its exact inputs
identify an irreducible actual double branched cover in each case.
An incompressible separating torus cuts off a certified hyperbolic
component, respectively m036, m121 or m189.
Thus the closed cover is not hyperbolic and is not a graph manifold.
Theorem~\ref{thm:class} gives the exclusion.
\end{proof}

The finite replay is \path{verify_rational_replacement_topology.py}.
Its input is \path{rational_replacement_topology_inputs.json}.
It enumerates all labelled four-vertex simple graphs, checks the
tetrahedral medial graph and its six four-ended sites, verifies the
Mayer--Vietoris dimension obstruction, and computes the mod-two
homology of all twenty-eight retained closed triangulations.
Those actual-cover bindings and incompressible-torus certificates
remain separate dependencies, sealed by their retained receipt;
they are not rerun by this smaller checker.
The new receipt is \path{rational_replacement_topology_verification.json}.

The quotient diagram in \cite{km} and the relevant classification
statement in \cite{mar} were read and visually checked in the
retained primary PDFs. The arborescent-cover construction, the complete
exceptional-filling classification and its computational ingredients
are classical imports, not independent replays here.
The local graph-to-tangle argument, the actual branched-cover gluing
and the mod-two exclusion are the written proof above.
A passing finite check is not formal verification or external
specialist review.

The exclusion concerns rational substitutions on a four-face base.
The next section shows why it cannot extend to arbitrary higher-face
origins when no rank or bridge preservation is required. Arbitrary
algebraic substitutions and other mechanisms proving MRC remain
outside Theorem~\ref{thm:class}.

\section{A concrete higher-face origin and the scope boundary}

\begin{theorem}\label{thm:eightbase}
There is a BBK twisted knot $J$ with $b(J)=\operatorname{mr}(J)=8$
and a displayed $35$-crossing diagram such that arbitrary rational
replacements of its twelve twist boxes give $13a_{373}$.
These replacements do not preserve the bridge number or the Coxeter
rank of the base. This is not a counterexample to the meridional rank
conjecture.
\end{theorem}
\begin{proof}
Start with the retained thirteen-crossing PD for $13a_{373}$, numbering
crossings from zero. Crossings $3,4$ are
\[
                 (13,4,14,5),\qquad(5,14,6,15).
\]
They form a two-crossing rational box with external cyclic ports
$(13,4,6,15)$. Replace this box by one crossing site, retaining its
boundary marking. The resulting twelve-site shadow is
\[
\begin{split}
 &(3,1,4,26),\ (1,7,2,6),\ (7,3,8,2),\ (13,4,6,15),\\
 &(21,9,22,8),\ (9,17,10,16),\ (17,11,18,10),\\
 &(11,25,12,24),\ (19,12,20,13),\ (15,22,16,23),\\
 &(23,18,24,19),\ (25,21,26,20).
\end{split}
\]
One of its checkerboard graphs $G$ has eight vertices and edge list
\[
 12,\ 01,\ 20,\ 17,\ 52,\ 35,\ 63,\ 46,\ 47,\ 57,\ 67,\ 24.
\]
Here $ab$ denotes the edge with endpoints $a,b$; the shadow specifies
its plane embedding. This graph is connected, loopless and simple.

Use $G$ as the dual of the fat-vertex graph $\Gamma$.
Put a twist band of absolute weight two at site $1$ and absolute
weight three at each of the other eleven sites. The band direction is
fixed by this duality. A band subdivides its edge in the full
crossing graph of $\Gamma$; dually, the full $G$ crossing graph
replaces each edge by its indicated number of parallel crossing edges.
For an exact local convention, number the four ports of a row
$(a,b,c,d)$ cyclically from zero. Let $p\in\{0,1\}$ be the index of
one of its $G$-coloured corners, a corner $h$ lying between ports
$h$ and $h+1$. One extra half-twist is inserted by
\[
\begin{array}{c|c}
p=0&(a,x,y,d),\ (x,b,c,y)\\
p=1&(a,b,x,y),\ (y,x,c,d),
\end{array}
\]
with fresh edge labels $x,y$. Repetition at the first row gives the
specified two- or three-crossing twist. The new bigon has the
$\Gamma$ colour, so each insertion adds a parallel edge in $G$.
For the displayed shadow the corner parities are
\[
                  (1,1,1,1,1,1,1,1,1,1,1,1).
\]
The exact input retains the resulting full $35$-crossing PD.
The finite checker reconstructs all these rows, verifies a spherical
connected rotation system and a single link component. It checks
the twelve parallel bundles on eight vertices and reduces them to the
simple dual $G$. The chosen crossing signs are uniform for this
checkerboard colouring, so the full diagram is alternating.
The full $G$ crossing graph is loopless, and has no bridges since
every bundle has at least two edges. Hence neither checkerboard
graph has a loop, excluding nugatory crossings and in particular
Reidemeister~I moves. Alternation excludes Reidemeister~II moves.
Thus the diagram is reduced as required by the definition, and
this is a twisted knot in the precise sense of
BBK Section~4: all band weights have absolute value at least two,
and $\Gamma^*=G$ is simple. Their Theorem~1 gives
$b(J)=\operatorname{mr}(J)=8$.

Replace the rational twist box at site $3$ by the original two-crossing
box displayed above, and replace every other box by its original
single crossing with its original ports. This restores the retained
PD literally. Every inserted box is rational. Since $13a_{373}$ has
bridge number four by its retained certificate, the operation has
dropped the bridge number from eight to four; the sharp lower and
upper bounds of the base are not preserved.
\end{proof}

The exact files for this finite example are
\path{higher_face_rational_origin_inputs.json},
\path{verify_higher_face_rational_origin.py} and
\path{higher_face_rational_origin_verification.json}.
The checker uses the standard library. It independently traverses
corner cycles, checks the parallel twist bundles, alternation, absence of loops and bridges, and strand connectivity, binds
the target PD to the retained named-knot input, and rejects changes
to the base or recovered diagram and a forbidden unit base weight.
The interpretation of a twist band and the BBK bridge theorem are
written dependencies; these finite checks are not formal verification.

\begin{theorem}\label{thm:universalorigin}
For every knot $K$ and every integer $N$, some BBK twisted knot $J$
with $b(J)\ge N$ has arbitrary rational replacements of its twist
boxes equal to $K$.
\end{theorem}
\begin{proof}
Choose a connected nonempty diagram of $K$, allowing a nugatory
crossing when $K$ is the unknot. Let $H$ be either checkerboard plane
graph, with its inherited cyclic orders. Subdivide every edge into a
path of length $\ell\ge3$, using different internal vertices for
different edges. The resulting plane graph $G$ is simple: a loop of
$H$ becomes a cycle of length at least three, and every edge of $G$
contains an internal vertex private to its original edge. Distinct
original edges therefore create no parallel pair. Increase $\ell$
until $\#V(G)\ge N$.

The medial box construction for $G$ is the medial construction for $H$
with extra degree-two vertices along its edges. A subdivision glues
two boxes by the adjacent-port Conway operation already described.
In one box on each original edge put its original signed single
crossing, and in every other box on that path put the identity for
that operation: the zero tangle for addition, or its quarter-turn
for the corresponding product. These are rational tangles. Removing
the identity boxes restores the original diagram with the same
four external ports on each original edge. Thus this rational
decoration of $G$ is exactly $K$.

It remains to choose a twisted \emph{knot} on the same base $G$.
Let $\Gamma=G^*$ and put two half-twists in every band.
Choose their handedness so that every crossing has the same
checkerboard sign. The resulting diagram is alternating.
The full crossing graph of the dual $G$ replaces every edge by
two parallel edges. Since $G$ is loopless and each edge has at least
two parallel copies, this crossing graph is loopless and bridge-free.
Thus neither checkerboard graph has a loop, so there is no nugatory
crossing or Reidemeister~I move. Alternation excludes a
Reidemeister~II move. This verifies the reduced-diagram hypothesis
in the BBK definition.

At a box, changing its weight between two and three, while keeping
the same checkerboard sign, changes its endpoint matching to a
different perfect matching of the same four ports. The resulting diagram is a union
of circuits in the connected medial shadow. If there is more than
one circuit, some site is traversed by two different circuits:
otherwise all four ports of each vertex would lie on one circuit,
and the connected shadow would partition into disjoint vertex sets.
At such a site, switch the weight from two to three or from three
to two. Either alternative matching reconnects the two distinct
circuits into a single circuit. No other circuit changes. Repeating
this step strictly reduces the number of circuits until there is
one. All band weights are still two or three in absolute value.
The checkerboard signs remain uniform, and the full dual crossing
graph still consists of parallel bundles of size at least two.
Thus alternation and the absence of Reidemeister~I and~II moves
hold throughout the circuit merges.
Consequently the result is a BBK twisted knot $J$, and their
Theorem~1 gives $b(J)=\#V(G)\ge N$.

Finally replace its boxes by the rational decoration recovering $K$
constructed in the first part. The base and its plane ports have not
changed, so these are the permitted arbitrary rational replacements.
\end{proof}

Theorem~\ref{thm:universalorigin} concerns geometric origins under
arbitrary rational replacement. It does not extend the rank-preserving
equality proof of BBK: a construction preserving a Coxeter
rank $r$ and a matching bridge upper bound still certifies that same
$r$. In particular, a sharp-four application of that proof uses
$r=4$. Starting from a higher-face base and allowing those bounds
to drop supplies no prior sharp-four certificate for the target.

\begin{thebibliography}{9}
\bibitem{bbk}
S. Baader, R. Blair and A. Kjuchukova,
Coxeter groups and meridional rank of links,
Math. Ann. 379 (2021), 1533--1551.
Theorem~1, Section~4 definition and subsequent replacement paragraph.
\url{https://doi.org/10.1007/s00208-020-02124-z}.
\bibitem{bs}
F. Bonahon and L. C. Siebenmann,
New Geometric Splittings of Classical Knots and the Classification
and Symmetries of Arborescent Knots, author monograph,
version 26 January 2016. Introduction pp.~vi--vii and Part~V.
\url{https://dornsife.usc.edu/francis-bonahon/wp-content/uploads/sites/205/2023/06/BonSieb-compressed.pdf}.
\bibitem{km}
A. Kolpakov and B. Martelli,
Hyperbolic four-manifolds with one cusp,
Geom. Funct. Anal. 23 (2013), 1904--1933.
Author preprint arXiv:1303.6122v7,
Figure~2 and Section~1.1, printed pp.~5--6.
\url{https://doi.org/10.1007/s00039-013-0247-2}.
\url{https://arxiv.org/abs/1303.6122v7}.
\bibitem{mar}
B. Martelli, Dehn surgery on the minimally twisted seven-chain link,
Comm. Anal. Geom. 29 (2021), 1597--1641.
Author preprint arXiv:1808.08430v2, Theorem~2.7, printed pp.~15--16,
and the explanation of its two mixed pieces.
\url{https://doi.org/10.4310/CAG.2021.v29.n7.a5}.
\url{https://arxiv.org/abs/1808.08430v2}.
\end{thebibliography}
\end{document}
